\documentclass[10pt,a4paper]{amsart}
\usepackage{amsmath,amssymb,amsthm,graphicx}
\usepackage{verbatim}
\newenvironment{reference}{\comment}{\endcomment}
%\newenvironment{reference}{}{}
\newenvironment{slogan}{\comment}{\endcomment}
\newenvironment{history}{\comment}{\endcomment}

% For commutative diagrams we use Xy-pic
\usepackage[all]{xy}

% We use 2cell for 2-commutative diagrams.
\xyoption{2cell}
% Theorem environments.
%
\theoremstyle{plain}
\newtheorem{theorem}[subsection]{Theorem}
\newtheorem{proposition}[subsection]{Proposition}
\newtheorem{lemma}[subsection]{Lemma}

\theoremstyle{definition}
\newtheorem{definition}[subsection]{Definition}
\newtheorem{example}[subsection]{Example}
\newtheorem{exercise}[subsection]{Exercise}
\newtheorem{situation}[subsection]{Situation}

\theoremstyle{remark}
\newtheorem{remark}[subsection]{Remark}
\newtheorem{remarks}[subsection]{Remarks}

\numberwithin{equation}{subsection}

% Macros
%
\def\lim{\mathop{\mathrm{lim}}\nolimits}
\def\colim{\mathop{\mathrm{colim}}\nolimits}
\def\Spec{\mathop{\mathrm{Spec}}}
\def\Hom{\mathop{\mathrm{Hom}}\nolimits}
\def\Ext{\mathop{\mathrm{Ext}}\nolimits}
\def\SheafHom{\mathop{\mathcal{H}\!\mathit{om}}\nolimits}
\def\SheafExt{\mathop{\mathcal{E}\!\mathit{xt}}\nolimits}
\def\Sch{\mathit{Sch}}
\def\Mor{\mathop{\mathrm{Mor}}\nolimits}
\def\Ob{\mathop{\mathrm{Ob}}\nolimits}
\def\Sh{\mathop{\mathit{Sh}}\nolimits}
\def\NL{\mathop{N\!L}\nolimits}
\def\CH{\mathop{\mathrm{CH}}\nolimits}
\def\proetale{{pro\text{-}\acute{e}tale}}
\def\etale{{\acute{e}tale}}
\def\QCoh{\mathit{QCoh}}
\def\Ker{\mathop{\mathrm{Ker}}}
\def\Im{\mathop{\mathrm{Im}}}
\def\Coker{\mathop{\mathrm{Coker}}}
\def\Coim{\mathop{\mathrm{Coim}}}

% Boxtimes
%
\DeclareMathSymbol{\boxtimes}{\mathbin}{AMSa}{"02}

%
% Macros for moduli stacks/spaces
%
\def\QCohstack{\mathcal{QC}\!\mathit{oh}}
\def\Cohstack{\mathcal{C}\!\mathit{oh}}
\def\Spacesstack{\mathcal{S}\!\mathit{paces}}
\def\Quotfunctor{\mathrm{Quot}}
\def\Hilbfunctor{\mathrm{Hilb}}
\def\Curvesstack{\mathcal{C}\!\mathit{urves}}
\def\Polarizedstack{\mathcal{P}\!\mathit{olarized}}
\def\Complexesstack{\mathcal{C}\!\mathit{omplexes}}
% \Pic is the operator that assigns to X its picard group, usage \Pic(X)
% \Picardstack_{X/B} denotes the Picard stack of X over B
% \Picardfunctor_{X/B} denotes the Picard functor of X over B
\def\Pic{\mathop{\mathrm{Pic}}\nolimits}
\def\Picardstack{\mathcal{P}\!\mathit{ic}}
\def\Picardfunctor{\mathrm{Pic}}
\def\Deformationcategory{\mathcal{D}\!\mathit{ef}}
\title{A local ring with a regular sequence of length three whose quotient has a nonzero element killed by two entries}
\author{The Stacks Project Authors}
\date{stacks-project:examples.tex:local-snapshot-20261012; GNU FDL 1.2 or later}
\begin{document}
\maketitle
\noindent\textit{Editorial note.} All delivered mathematical text is the author's own, copied from the pinned Stacks Project chapter file examples.tex; no statement or proof step was rewritten, reordered or removed.\par\medskip
\noindent\textit{Editorial note.} The counted proof is the author's complete argument: Hironaka-style example; the companion lemma lemma-base-change-regular-sequence consumes this construction.\par\medskip
\noindent\textit{Editorial note.} The label lemma-strange-regular-sequence names the counted statement, and the proof environment immediately following it in the source is delivered verbatim as its proof.\par\medskip
\noindent\textit{Editorial note.} The counted statement is carried by the author's own lemma environment in the chapter source.\par\medskip
\noindent\textit{Allowed standard prerequisites (editorial).}\begin{itemize}
\item Section construction of the module E: E is the module of equivalence classes of pairs (f, g) with f in k[x, y, z, x inverse, y inverse] and g in k[x, y, z, y inverse], under the equivalence generated by f equals f-prime plus x h modulo y times k[x, y, z, x inverse] and g equals g-prime minus z h modulo y z times k[x, y, z]; the section prose proves (a) multiplication by x is injective on E, (b) multiplication by y is injective on E modulo xE, (c) E over the ideal generated by x and y is zero, and (d) the class of the pair (1, 0) yields a nonzero element delta in E over zE with x delta = y delta = 0
\item Square-zero extension: for a ring S and an ideal I with I squared zero, the ring S plus I has maximal ideal the sum of the maximal ideal of S and I whenever that sum is an ideal
\item Localization is exact, so a regular sequence remains regular after localizing at a multiplicative set
\end{itemize}\medskip
\begin{lemma}
\label{lemma-strange-regular-sequence}
There exists a local ring $R$ and a regular sequence $x, y, z$
(in the maximal ideal) such that there exists a nonzero element
$\delta \in R/zR$ with $x\delta = y\delta = 0$.
\end{lemma}

\begin{proof}
Let $R = k[x, y, z] \oplus E$ where $E$ is the module above considered
as a square zero ideal. Then it is clear that $x, y, z$ is a regular
sequence in $R$, and that the element $\delta \in E/zE \subset R/zR$
gives an element with the desired properties. To get a local example
we may localize $R$ at the maximal ideal $\mathfrak m = (x, y, z, E)$.
The sequence $x, y, z$ remains a regular sequence (as localization is
exact), and the element $\delta$ remains nonzero as it is supported
at $\mathfrak m$.
\end{proof}

\end{document}
